\documentclass[10pt,a4paper]{amsart}
\usepackage{amsmath,amssymb,amsthm,graphicx}
\usepackage{verbatim}
\newenvironment{reference}{\comment}{\endcomment}
%\newenvironment{reference}{}{}
\newenvironment{slogan}{\comment}{\endcomment}
\newenvironment{history}{\comment}{\endcomment}

% For commutative diagrams we use Xy-pic
\usepackage[all]{xy}

% We use 2cell for 2-commutative diagrams.
\xyoption{2cell}
% Theorem environments.
%
\theoremstyle{plain}
\newtheorem{theorem}[subsection]{Theorem}
\newtheorem{proposition}[subsection]{Proposition}
\newtheorem{lemma}[subsection]{Lemma}

\theoremstyle{definition}
\newtheorem{definition}[subsection]{Definition}
\newtheorem{example}[subsection]{Example}
\newtheorem{exercise}[subsection]{Exercise}
\newtheorem{situation}[subsection]{Situation}

\theoremstyle{remark}
\newtheorem{remark}[subsection]{Remark}
\newtheorem{remarks}[subsection]{Remarks}

\numberwithin{equation}{subsection}

% Macros
%
\def\lim{\mathop{\mathrm{lim}}\nolimits}
\def\colim{\mathop{\mathrm{colim}}\nolimits}
\def\Spec{\mathop{\mathrm{Spec}}}
\def\Hom{\mathop{\mathrm{Hom}}\nolimits}
\def\Ext{\mathop{\mathrm{Ext}}\nolimits}
\def\SheafHom{\mathop{\mathcal{H}\!\mathit{om}}\nolimits}
\def\SheafExt{\mathop{\mathcal{E}\!\mathit{xt}}\nolimits}
\def\Sch{\mathit{Sch}}
\def\Mor{\mathop{\mathrm{Mor}}\nolimits}
\def\Ob{\mathop{\mathrm{Ob}}\nolimits}
\def\Sh{\mathop{\mathit{Sh}}\nolimits}
\def\NL{\mathop{N\!L}\nolimits}
\def\CH{\mathop{\mathrm{CH}}\nolimits}
\def\proetale{{pro\text{-}\acute{e}tale}}
\def\etale{{\acute{e}tale}}
\def\QCoh{\mathit{QCoh}}
\def\Ker{\mathop{\mathrm{Ker}}}
\def\Im{\mathop{\mathrm{Im}}}
\def\Coker{\mathop{\mathrm{Coker}}}
\def\Coim{\mathop{\mathrm{Coim}}}

% Boxtimes
%
\DeclareMathSymbol{\boxtimes}{\mathbin}{AMSa}{"02}

%
% Macros for moduli stacks/spaces
%
\def\QCohstack{\mathcal{QC}\!\mathit{oh}}
\def\Cohstack{\mathcal{C}\!\mathit{oh}}
\def\Spacesstack{\mathcal{S}\!\mathit{paces}}
\def\Quotfunctor{\mathrm{Quot}}
\def\Hilbfunctor{\mathrm{Hilb}}
\def\Curvesstack{\mathcal{C}\!\mathit{urves}}
\def\Polarizedstack{\mathcal{P}\!\mathit{olarized}}
\def\Complexesstack{\mathcal{C}\!\mathit{omplexes}}
% \Pic is the operator that assigns to X its picard group, usage \Pic(X)
% \Picardstack_{X/B} denotes the Picard stack of X over B
% \Picardfunctor_{X/B} denotes the Picard functor of X over B
\def\Pic{\mathop{\mathrm{Pic}}\nolimits}
\def\Picardstack{\mathcal{P}\!\mathit{ic}}
\def\Picardfunctor{\mathrm{Pic}}
\def\Deformationcategory{\mathcal{D}\!\mathit{ef}}
\usepackage{cleveref}
\crefname{originallemma}{original lemma}{original lemmas}
\Crefname{originallemma}{Original lemma}{Original lemmas}
\makeatletter
\newlabel{algebra-lemma-characterize-reduced-ring-normal}{{030C}{source}{}{sourcepointer-1}{}}
\newlabel{algebra-lemma-characterize-reduced-ring-normal@cref}{{[originallemma][0][]030C}{[1][0][]source}{}{}{}}
\makeatother
\crefname{originaldefinition}{original definition}{original definitions}
\Crefname{originaldefinition}{Original definition}{Original definitions}
\makeatletter
\newlabel{decent-spaces-definition-very-reasonable}{{03I8}{source}{}{sourcepointer-2}{}}
\newlabel{decent-spaces-definition-very-reasonable@cref}{{[originaldefinition][0][]03I8}{[1][0][]source}{}{}{}}
\makeatother
\crefname{originallemma}{original lemma}{original lemmas}
\Crefname{originallemma}{Original lemma}{Original lemmas}
\makeatletter
\newlabel{decent-spaces-lemma-decent-generic-points}{{0ABV}{source}{}{sourcepointer-3}{}}
\newlabel{decent-spaces-lemma-decent-generic-points@cref}{{[originallemma][0][]0ABV}{[1][0][]source}{}{}{}}
\makeatother
\title{A normal integral algebraic space is covered by normal integral affine schemes}
\author{The Stacks Project Authors}
\date{stacks-project:spaces-over-fields.tex:local-snapshot-20261012; GNU FDL 1.2 or later}
\begin{document}
\maketitle
\noindent\textit{Editorial note.} All delivered mathematical text is the author's own, copied from the pinned Stacks Project chapter file spaces-over-fields.tex; no statement or proof step was rewritten, reordered or removed.\par\medskip
\noindent\textit{Editorial note.} The counted proof is the author's complete argument: the finiteness of the number of irreducible components of an affine etale neighbourhood through decency, the decomposition of the normal ring of the neighbourhood into a product of normal domains, and the selection of the component through which the given point is reached.\par\medskip
\noindent\textit{Editorial note.} The label lemma-normal-integral-cover-by-affines names the counted statement, and the proof environment immediately following it in the source is delivered verbatim as its proof.\par\medskip
\noindent\textit{Editorial note.} The counted statement is carried by the author's own lemma environment in the chapter source.\par\medskip
\noindent\textit{Editorial note.} Pre-open gate: python operations/agent-authoring/tools/precard\_gate.py --source spaces-over-fields.tex --label lemma-normal-integral-cover-by-affines --start 516 --end 536 --local-id 28067 returned RESULT=OK (L28067 free, occupancy FREE, public interval clear). Candidate taken from sources/stacks-spaces-over-fields/candidates-halfA.json.\par\medskip
\noindent\textit{Editorial note.} The statement and proof are copied from the pinned chapter file; the cited inputs are kept as external pointers with their canonical Stacks tags.\par\medskip
\noindent\textit{Allowed standard prerequisites (editorial).}\begin{itemize}
\item Finiteness of the number of irreducible components of an etale neighbourhood. Decent Spaces, the definition of a decent space and the lemma that the generic points of an etale neighbourhood map to the generic point of the space. Supplies the finiteness of the number of generic points of the affine neighbourhood, hence of its minimal primes.
\item Decomposition of a normal ring with finitely many minimal primes. Algebra, the lemma characterizing a reduced ring as a product of domains indexed by its minimal primes, applied to a normal ring. Supplies the decomposition of the ring of the neighbourhood into a product of normal domains.
\end{itemize}\medskip
\ifdefined\hypertarget\hypertarget{sourcepointer-1}{}\fi
\noindent\textit{Source reference.} Original lemma 030C: sources/stacks-chow/tags.txt, source lines 3943--3943. Prerequisite pointer invoked at proof line(s) 533; it names a standard fact of a sibling chapter that the counted argument uses, and every step of the counted proof built on it is retained verbatim in the delivered body.\par\medskip
\ifdefined\hypertarget\hypertarget{sourcepointer-2}{}\fi
\noindent\textit{Source reference.} Original definition 03I8: sources/stacks-chow/tags.txt, source lines 4594--4594. Prerequisite pointer invoked at proof line(s) 529; it names a standard fact of a sibling chapter that the counted argument uses, and every step of the counted proof built on it is retained verbatim in the delivered body.\par\medskip
\ifdefined\hypertarget\hypertarget{sourcepointer-3}{}\fi
\noindent\textit{Source reference.} Original lemma 0ABV: sources/stacks-chow/tags.txt, source lines 13033--13033. Prerequisite pointer invoked at proof line(s) 527; it names a standard fact of a sibling chapter that the counted argument uses, and every step of the counted proof built on it is retained verbatim in the delivered body.\par\medskip
\begin{lemma}
\label{lemma-normal-integral-cover-by-affines}
Let $S$ be a scheme. Let $X$ be a normal integral algebraic space over $S$.
For every $x \in |X|$ there exists a normal integral affine scheme $U$
and an \'etale morphism $U \to X$ such that $x$ is in the image.
\end{lemma}

\begin{proof}
Choose an affine scheme $U$ and an \'etale morphism $U \to X$ such that
$x$ is in the image. Let $u_i$, $i \in I$ be the generic points of irreducible
components of $U$. Then each $u_i$ maps to the generic point of $X$
(Decent Spaces, Lemma \ref{decent-spaces-lemma-decent-generic-points}). By 
our definition of a decent space
(Decent Spaces, Definition \ref{decent-spaces-definition-very-reasonable}),
we see that $I$ is finite. Hence $U = \Spec(A)$ where $A$ is a normal ring
with finitely many minimal primes.
Thus $A = \prod_{i \in I} A_i$ is a product of normal domains by
Algebra, Lemma \ref{algebra-lemma-characterize-reduced-ring-normal}.
Then $U = \coprod U_i$ with $U_i = \Spec(A_i)$ and $x$ is in the image of
$U_i \to X$ for some $i$. This proves the lemma.
\end{proof}

\end{document}
